\section{Simulation Experimental Details}~\label{supp:sec:simulation_setup_details}
In this section, we provide additional details of our simulation experiments, including: (1) the training tasks in our setup (Sec.~\ref{supp:sec:subsec:training_tasks}), (2) how we obtain 2D grounding information in simulation (Sec.~\ref{supp:sec:subsec:grounding-source}), (3) the per-model training configurations (Sec.~\ref{supp:sec:subsec:sim-training-details}), and (4) the full experimental results (Sec.~\ref{supp:sec:subsec:sim-experimental-results-details}).

\subsection{Training Data Details}
\label{supp:sec:subsec:training_tasks}
Our default training setup augments LIBERO-Object and LIBERO-Spatial with 49 additional pick-and-place tasks from LIBERO-90, yielding 69 tasks in total. 
In Tab.~\ref{tab:libero90_tasks}, we list the 49 additional pick-and-place tasks.
Our selection criterion is that a mask for the target object or region must be obtainable from the observation. For abstract regions without an explicit mask (e.g., the left side of a mug), we instead ensure that the 2D projection of the region's center is not occluded. For example, tasks with target region such as the second shelf of a cabinet are excluded from our selection, since their center, when projected to 2D, falls on the surface of the top shelf and is therefore occluded.
\begin{table}[t]
\centering
\footnotesize
\caption{The 49 tasks from LIBERO-90 used in the 69-task setup.}
\label{tab:libero90_tasks}
\begin{tabular}{ll}
\toprule
Scene & Task Instruction \\
\midrule
\texttt{KITCHEN\_SCENE1}      & put the black bowl on the plate \\
\texttt{KITCHEN\_SCENE1}      & put the black bowl on top of the cabinet \\
\texttt{KITCHEN\_SCENE2}      & put the black bowl at the back on the plate \\
\texttt{KITCHEN\_SCENE2}      & put the black bowl at the front on the plate \\
\texttt{KITCHEN\_SCENE2}      & put the middle black bowl on the plate \\
\texttt{KITCHEN\_SCENE2}      & put the middle black bowl on top of the cabinet \\
\texttt{KITCHEN\_SCENE2}      & stack the black bowl at the front on the black bowl in the middle \\
\texttt{KITCHEN\_SCENE2}      & stack the middle black bowl on the back black bowl \\
\texttt{KITCHEN\_SCENE3}      & put the frying pan on the stove \\
\texttt{KITCHEN\_SCENE3}      & put the moka pot on the stove \\
\texttt{KITCHEN\_SCENE4}      & put the black bowl in the bottom drawer of the cabinet \\
\texttt{KITCHEN\_SCENE4}      & put the black bowl on top of the cabinet \\
\texttt{KITCHEN\_SCENE4}      & put the wine bottle in the bottom drawer of the cabinet \\
\texttt{KITCHEN\_SCENE5}      & put the black bowl in the top drawer of the cabinet \\
\texttt{KITCHEN\_SCENE5}      & put the black bowl on the plate \\
\texttt{KITCHEN\_SCENE5}      & put the black bowl on top of the cabinet \\
\texttt{KITCHEN\_SCENE5}      & put the ketchup in the top drawer of the cabinet \\
\texttt{KITCHEN\_SCENE6}      & put the yellow and white mug to the front of the white mug \\
\texttt{KITCHEN\_SCENE7}      & put the white bowl on the plate \\
\texttt{KITCHEN\_SCENE7}      & put the white bowl to the right of the plate \\
\texttt{KITCHEN\_SCENE8}      & put the right moka pot on the stove \\
\texttt{KITCHEN\_SCENE9}      & put the frying pan on top of the cabinet \\
\texttt{KITCHEN\_SCENE9}      & put the white bowl on top of the cabinet \\
\texttt{KITCHEN\_SCENE10}     & put the black bowl in the top drawer of the cabinet \\
\texttt{LIVING\_ROOM\_SCENE1} & pick up the alphabet soup and put it in the basket \\
\texttt{LIVING\_ROOM\_SCENE1} & pick up the cream cheese box and put it in the basket \\
\texttt{LIVING\_ROOM\_SCENE1} & pick up the ketchup and put it in the basket \\
\texttt{LIVING\_ROOM\_SCENE1} & pick up the tomato sauce and put it in the basket \\
\texttt{LIVING\_ROOM\_SCENE2} & pick up the butter and put it in the basket \\
\texttt{LIVING\_ROOM\_SCENE2} & pick up the milk and put it in the basket \\
\texttt{LIVING\_ROOM\_SCENE2} & pick up the orange juice and put it in the basket \\
\texttt{LIVING\_ROOM\_SCENE2} & pick up the tomato sauce and put it in the basket \\
\texttt{LIVING\_ROOM\_SCENE3} & pick up the alphabet soup and put it in the tray \\
\texttt{LIVING\_ROOM\_SCENE3} & pick up the butter and put it in the tray \\
\texttt{LIVING\_ROOM\_SCENE3} & pick up the cream cheese and put it in the tray \\
\texttt{LIVING\_ROOM\_SCENE3} & pick up the ketchup and put it in the tray \\
\texttt{LIVING\_ROOM\_SCENE3} & pick up the tomato sauce and put it in the tray \\
\texttt{LIVING\_ROOM\_SCENE4} & pick up the black bowl on the left and put it in the tray \\
\texttt{LIVING\_ROOM\_SCENE4} & pick up the salad dressing and put it in the tray \\
\texttt{LIVING\_ROOM\_SCENE4} & pick up the chocolate pudding and put it in the tray \\
\texttt{LIVING\_ROOM\_SCENE5} & put the red mug on the left plate \\
\texttt{LIVING\_ROOM\_SCENE5} & put the red mug on the right plate \\
\texttt{LIVING\_ROOM\_SCENE5} & put the white mug on the left plate \\
\texttt{LIVING\_ROOM\_SCENE5} & put the yellow and white mug on the right plate \\
\texttt{LIVING\_ROOM\_SCENE6} & put the chocolate pudding to the left of the plate \\
\texttt{LIVING\_ROOM\_SCENE6} & put the chocolate pudding to the right of the plate \\
\texttt{LIVING\_ROOM\_SCENE6} & put the red mug on the plate \\
\texttt{LIVING\_ROOM\_SCENE6} & put the white mug on the plate \\
\texttt{STUDY\_SCENE4}        & pick up the book on the left and place it on top of the shelf \\
\bottomrule
\end{tabular}
\end{table}

\subsection{Grounding Source Details}
\label{supp:sec:subsec:grounding-source}

We obtain 2D oracle target points from simulator-provided segmentation and object-center information, then lift them into 3D following Sec.~\ref{supp:sec:subsec:lifting}. A key invariant underlies our sampling: every 3D target point must be obtainable from a camera-visible 2D pixel via depth-based lifting, never taken directly as a 3D centroid from the simulator. This guarantees that our pipeline genuinely starts from a 2D grounding signal—the same signal an off-the-shelf VLM or detector could provide—rather than relying on privileged 3D information unavailable at deployment.

To respect this invariant, we adopt two sampling procedures depending on whether the sub-goal target admits an explicit segmentation mask. For object-typed targets, such as the target object or placement regions like \textit{plate} or \textit{tray}, the mask is directly available, and we apply a distance transform~\cite{distancetransform} to select the innermost pixel, i.e., the one with maximum distance to the mask boundary. For abstract placement regions that lack an explicit mask, such as \textit{to the left of the plate}, we instead start from the region's 3D centroid, project it onto the image to obtain a 2D pixel, and then re-lift this pixel into 3D via the same depth-based procedure. The resulting target point is the camera-observed surface point along the ray through the projected centroid, ensuring that even abstract regions are grounded through a visible 2D pixel rather than a free-space coordinate.

\subsection{Training Details}
\label{supp:sec:subsec:sim-training-details}
We detail the fine-tuning setups of GR00T-N1.6~\cite{gr00tn1_2025} and $\pi_{0.5}$~\cite{pi05} in simulation experiments below.

\noindent\textbf{Training Details for GR00T-N1.6.}
We fine-tune GR00T-N1.6 in simulation following the official training recipe on LIBERO dataset, with the only modification being the global batch size. The model predicts an action chunk of $16$ future steps per forward pass. We use a global batch size of $384$ and train for $20{,}000$ steps using AdamW with a peak learning rate of $1{\times}10^{-4}$, a cosine decay schedule with a warmup ratio of $0.05$ ($\sim\!1{,}000$ steps), weight decay of $1{\times}10^{-5}$, and gradient clipping at a max norm of $1.0$. For visual augmentation, input images are resized so that the shortest edge is $256$ px, randomly cropped to a square patch covering $0.95$ of the shortest edge, and then resized back to $256$ px. We further apply ColorJitter-style perturbations with brightness $0.3$, contrast $0.4$, saturation $0.5$, and hue $0.08$, while random rotation is disabled. All other hyperparameters follow the official configuration~\cite{gr00tn1_2025}.

\noindent\textbf{Training Details for $\pi_{0.5}$.}
We fine-tune $\pi_{0.5}$ in simulation following the official training recipe on LIBERO dataset. The model predicts an action chunk of $10$ future steps per forward pass. We use a global batch size of $256$ and train for $30{,}000$ steps using AdamW with a peak learning rate of $5{\times}10^{-5}$, a linear warmup schedule over the first $10{,}000$ steps and held constant thereafter, weight decay of $1{\times}10^{-10}$, and gradient clipping at a max norm of $1.0$. For visual augmentation, input images are resized (with padding) to $224{\times}224$ px; on the third-person view we apply a random crop covering $0.95$ of each spatial dimension, a bilinear resize back to $224{\times}224$ px, and a random rotation of $\pm 5^{\circ}$, while the wrist view receives no geometric augmentation. We further apply ColorJitter-style perturbations with brightness $0.3$, contrast $0.4$, and saturation $0.5$, while hue jitter is disabled. All other hyperparameters follow the official configuration~\cite{pi05}.


\subsection{Full Experimental Results}
\label{supp:sec:subsec:sim-experimental-results-details}
\definecolor{mygray}{HTML}{EFEFEF}
\begin{table}[t]
\centering
\caption{Controlled comparison of representation and injection on LIBERO-PRO under 20 training-task setup.}
\label{tab:controlled_experiment_20_tasks}
\resizebox{0.75\linewidth}{!}{
\begin{tabular}{l ccc ccc}
\toprule
\multirow{2}{*}{\textbf{Method}} & \multicolumn{3}{c}{Task Perturbation} & \multicolumn{3}{c}{Position Perturbation} \\
\cmidrule(lr){2-4} \cmidrule(lr){5-7}
  & Object & Spatial & Avg. & Object & Spatial & Avg. \\
\midrule

% Ours
GR00T-N1.6~\cite{gr00tn1_2025} & 9.0 & 48.8 & 28.9 & 4.8 & 2.0 & 3.4 \\
\midrule
+ 2D via Text Prompt & 10.0 & 54.2 & 32.1 & 24.6 & 6.0 & 15.3 \\
+ 2D via Visual Prompt & 10.0 & 49.2 & 29.6 & 1.2 & 2.2 & 1.7 \\
+ 2D via AdaLN & 11.2 & 68.0 & 39.6 & 35.4 & 3.4 & 19.4 \\
+ 3D via Text Prompt & 10.0 & 55.8 & 32.9 & 7.4 & 3.4 & 5.4 \\
+ 3D via AdaLN (Ours) & \textbf{52.4} & \textbf{71.6} & \textbf{62.0} & \textbf{67.8} & \textbf{15.2} & \textbf{41.5} \\
\bottomrule
\end{tabular}
}
\end{table}
\definecolor{mygray}{HTML}{EFEFEF}
\begin{table}[t]
\centering
\caption{Depth noise experiment on LIBERO-PRO under different training/testing depth noise condition.}
\label{tab:depth_noise_experiment_table}
\resizebox{0.75\linewidth}{!}{
\begin{tabular}{l cc cc c}
\toprule
\multirow{2}{*}{\textbf{(Training, Testing)}} & \multicolumn{2}{c}{Task Perturbation} & \multicolumn{2}{c}{Position Perturbation} & \multirow{2}{*}{\textbf{Average}} \\
\cmidrule(lr){2-3} \cmidrule(lr){4-5}
  & Object & Spatial & Object & Spatial & \\
\midrule
\textcolor{gray}{GR00T-N1.6 (Baseline)} & \textcolor{gray}{10.0} & \textcolor{gray}{52.4} & \textcolor{gray}{29.2} & \textcolor{gray}{27.0} & \textcolor{gray}{29.7} \\
\midrule
% Clean depth training
(Clean depth, Clean depth)        & 70.4 & 84.6 & 79.4 & 41.0 & 68.9 \\
(Clean depth, $\sigma=0.01$)      & 70.2 & 83.4 & 78.0 & 41.8 & 68.4 \\
(Clean depth, $\sigma=0.03$)      & 72.0 & 88.2 & 77.4 & 43.0 & 70.2 \\
\midrule
% sigma = 0.01 training
($\sigma=0.01$, Clean depth)      & 55.8 & 84.0 & 61.6 & 44.8 & 61.6 \\
($\sigma=0.01$, $\sigma=0.01$)    & 53.6 & 81.4 & 61.8 & 48.2 & 61.3 \\
($\sigma=0.01$, $\sigma=0.03$)    & 54.6 & 83.6 & 62.4 & 45.2 & 61.5 \\
\midrule
% sigma = 0.03 training
($\sigma=0.03$, Clean depth)      & 55.2 & 74.6 & 61.0 & 40.8 & 57.9 \\
($\sigma=0.03$, $\sigma=0.01$)    & 56.4 & 74.4 & 62.4 & 44.0 & 59.3 \\
($\sigma=0.03$, $\sigma=0.03$)    & 54.6 & 72.4 & 59.8 & 41.8 & 57.2 \\

\bottomrule
\end{tabular}
}
\end{table}
\definecolor{mygray}{HTML}{EFEFEF}
\begin{table}[t]
\centering
\caption{Grounding noise experiment on LIBERO-PRO under different grounding noise condition.}
\label{tab:grounding_noise_experiment_table}
\resizebox{0.75\linewidth}{!}{
\begin{tabular}{l cc cc c}
\toprule
\multirow{2}{*}{\textbf{Noise Condition}} & \multicolumn{2}{c}{Task Perturbation} & \multicolumn{2}{c}{Position Perturbation} & \multirow{2}{*}{\textbf{Average}} \\
\cmidrule(lr){2-3} \cmidrule(lr){4-5}
 & Object & Spatial & Object & Spatial & \\
\midrule
\textcolor{gray}{GR00T-N1.6 (Baseline)} & \textcolor{gray}{10.0} & \textcolor{gray}{52.4} & \textcolor{gray}{29.2} & \textcolor{gray}{27.0} & \textcolor{gray}{29.7} \\
\midrule
% Clean grounding noise testing
% \rowcolor{mygray} \multicolumn{6}{l}{\textit{Ours}} \\
$\sigma=0$ & 70.4 & 84.6 & 79.4 & 41.0 & 68.9 \\ 
$\sigma=10\%$ & 67.4 & 81.6 & 78.6 & 41.4 & 67.3 \\
$\sigma=30\%$ & 62.0 & 80.4 & 67.4 & 38.4 & 62.1 \\
$\sigma=50\%$ & 59.6 & 78.4 & 63.8 & 37.6 & 59.9 \\
Random sampling & 57.4 & 76.2 & 66.2 & 37.2 & 59.3 \\
\bottomrule
\end{tabular}
}
\end{table}
We report numerical performance values for all evaluations in Sec.~\ref{sec:tasks_scaling}, Sec.~\ref{sec:depth_noise_exp} and Sec.~\ref{sec:gnd_noise_exp}. Tab.~\ref{tab:controlled_experiment_20_tasks} presents the numerical results of the controlled comparison experiment on LIBERO-PRO under the 20-task setup (LIBERO-Object + LIBERO-Spatial), Tab.~\ref{tab:depth_noise_experiment_table} presents the numerical results of the depth noise experiment, and Tab.~\ref{tab:grounding_noise_experiment_table} presents the numerical results of the grounding noise experiment.